\scanpage{43}
% L1 — full passage crossed with long diagonal strokes
\deleted{si le \deleted{\unsure{fai}} motif de coefficient $\mathbb Z(0)$ est remplacé par \deleted{\unsure{le|un}} un motif constant tordu de \inserted{\unsure{purement|primitif}} type dimensionnel \deleted{\unsure{$d$}} $\delta$ [Attention au décalage par $\delta$ !]. Cet énoncé \deleted{\unsure{est}} \deleted{\unsure{de}} \uline{précise} les interprétations \unsure{galoisiennes|folkloriques|géométriques} et la conjecture de Tate du \deleted{\unsure{n°}}\dots, et est sans doute moins profond.}

% note crossing the boundary of the two sections
\margin{\deleted{$X$ de dim $n$, $U$ ouvert, $X-U$ de dim $<n$. Alors si $i\leq n$, $H^i(X)\deleted{\unsure{\simeq}}\longrightarrow H^i(U)$ est un isom \uuline{mod} $D_{i-1}(M^+_{\deleted{\unsure{i}}}(K))$, donc $H^i(X)$ et $H^i(U)$ ont $=$ composant ds $D_{i}(M^+(K))/D_{\unsure{n-1|i-1}}(M^+(K))$. De plus (via résolution) ce dernier est dans l'im de $M^+_{\unsure{i}}(K)$, d'où un inv. \unsure{birationel|birat.|birationnel} $\in\mathbb Z^{\sigma_i}$ \dots}}

% L9 — new heading
12) \uline{\deleted{\unsure{Bis|B.}}}\quad \uline{L'hom fondamental}
\[
L(K)\mathrel{\unsure{-|\longrightarrow}}M^+(K)
\]
\uline{et invariants birationnels fondamentaux}.

% L12
$K$ est un corps. On désigne par $L(K)$ le gpe $K$ engendré par les préschémas de type fini sur $K$, avec comme relations celles définies par découpage en morceaux. C'est un anneau par $\otimes$. On définit un hom d'anneaux
\[
\sigma:\quad L(K)\longrightarrow M^+(K)
\]
% archive sheet number: 42
